\centerline{\Huge \bf Индукция в графах}

\begin{flushright}
    \scriptsize{}
\end{flushright}

\textbf{Важное замечание.} \textit{ Предположим, мы решаем задачу про графы индукцией по числу вершин. Плохо говорить ''Возьмём граф на $n$ вершинах и добавим ещё одну вершину''. Причина здесь в следующем. Рассмотрим множества $A_n$ и $A_{n+1}$ графов на $n$ и $n+1$ вершине соответственно. Добавляя вершину (и какие-то рёбра) мы строим по графу из $A_n$ граф из $A_{n+1}$. Но мы необязательно сможем получить всевозможные графы из $A_{n+1}$. Правильно говорить ''Возьмём граф на $n+1$ вершине, удовлетворяющий условию, и выкинем из него одну вершину, со всеми идущими из неё рёбрами. К оставшемуся графу применим предположение индукции''.}

\problem{1} Докажите, что существует граф с $2n$ вершинами, степени которых равны $1, 1, 2, 2, \ldots,\\ n, n$.

\textbf{Решение.} База для $k = 1$ очевидна. Пусть для $k$ граф с $2k$ вершинами существует, докажем для $k+1$. Возьмём граф с $2k$ вершинами и разобьём его вершины на две части так, чтобы в каждой части были вершины степени $1, 2, \ldots, n$. Добавим две новые вершины, одну из них соединим со всеми вершинами  одной из частей, а также с еще одной новой. Несложно понять, что полученный граф удовлетворяет условию.

\problem{2} В стране любые два города соединены дорогой с односторонним движением. Докажите, что стартуя из некоторого города, можно проехать по всем городам, побывав в каждом по одному разу.
    
\textbf{Решение.} Докажем утверждение задачи индукцией по количеству вершин. База для 2 городов очевидна. Предположим, что утверждение верно для $n$ городов, докажем для $n+1$ города. Выкинем один из городов $B$ и все дороги, с концом в этом городе. По предположению индукции в оставшемся графе существует путь $A_1, A_2, \ldots, A_n$ по всем городам. Вернём город $B$. Если дорога ведёт из города $B$ в $A_1$, то начав с $B$, можно пройти обойти все города. Заметим, что если для некоторого $k$ из города $A_k$ дорога ведёт в $B$, а из $B$ дорога ведёт в $A_{k+1}$, то все города можно обойти следующим образом: $A_1 \to \ldots \to A_k \to B \to A_{k+1} \to \ldots \to A_n$. Если такого $k$ не существует, то в город $B$ ведёт дорога из каждого из городов $A_1, \ldots, A_n$, поэтому города можно обойти следующим образом: $A_1 \to \ldots \to A_n \to B$.

\newpage

\hspace{3mm}

\problem{3} Усадьбы любых двух графов в Истринском районе соединены либо водным, либо сухопутным сообщением.Докажите, что можно закрыть один из видов транспорта так, чтобы любой граф мог по-прежнему добраться до любого другого.
    
\textbf{Решение.} Для двух городов утверждение очевидно. Пусть верно для $n$ городов, докажем для $n+1$. Пусть сухопутные дороги --- это рёбра красного цвета, а водные --- синего. Выкинем одну вершину $B$ и все рёбра, ведущие из неё. Для остальных городов верно предположение индукции, то есть можно оставить только рёбра какого-то цвета (без ограничения общности синего) так, что граф будет связным. Вернём вершину $B$. Если из вершины $B$ выходит хотя бы одно ребро синего цвета, то если оставить только рёбра синего цвета, то граф по-прежнему будет связным. Если же из вершины $B$ выходят только рёбра красного цвета, то уберём все рёбра синего цвета. Оставшийся граф будет связным, поскольку от любой вершины можно добраться до любой через вершину $B$.

\problem{4} В некоторой стране каждый город соединен с каждым дорогой с односторонним движением. Докажите, что найдется город, из которого можно добраться в любой другой не более чем с одной пересадкой.
    
\textbf{Решение.} Для двух городов утверждение очевидно. Пусть верно для $n$ городов, докажем для $n+1$. Выкинем некоторую вершину $B$ и все рёбра, с концом в этой вершине. Для оставшегося графа верно предположение индукции, то есть существует вершина $A$ такая, что из него можно добраться до любой другой не более, чем 2 рёбрам. Пусть из вершины $A$ напрямую можно добраться в вершины $C_1, \ldots, C_k$, а только через другую вершину --- в $D_1, \ldots, D_l$. Если из вершины $A$ или из какой-нибудь из вершин $C_1, \ldots, C_k$ ребро ведёт в $B$, то вершина $A$ по-прежнему подходит. Если же во все эти вершины ведут рёбра из $B$, то подойдёт вершина $B$: в вершины $A, C_1, ldots, C_k$ можно попасть непосредственно, а в вершины $D_1, \ldots, D_l$ через вершины $C_1, \ldots, C_k$.

\problem{5} В компании из $n$ человек среди любых четверых есть знакомый с остальными тремя. Докажите, что есть человек, который знает всех остальных.
    
\textbf{Решение.} База для 4 человек очевидна. Пусть утверждение задачи верно для $n$ человек, докажем для $n+1$. Выкинем одну вершину $B$ и все рёбра, ведущие из неё. 

\newpage

\hspace{3mm}

\noindentДля оставшихся $n$ по предположению индукция найдётся вершина $A$, соединённая со всеми остальными. Если все остальные вершины также соединены между собой, то утверждение задачи очевидно --- вершина, которая соединена с $B$, будет искомой. Если же найдутся две вершины $C$ и $D$, не соединённые ребром, то рассмотрим четверку $A,B,C,D$. Заметим, что $C$ и $D$ не могут быть знакомы с $B$ одновременно, поскольку у них нет ребра между собой. Поэтому это будут одно и та же вершина $B$, но тогда эти вершина соединена ребром, то есть $A$ --- искомая вершина.

